\begin{lemma}[Simultaneous unequal retention]
Let $G$ be a simple graph on a finite vertex set, let $U,W$ be disjoint
sets of size $n>0$, and let $h>0$. Let $p:V\times V\to[0,1]$ be zero
on nonedges. Assume
\[
 2\,4^n\exp(-h^4n^2/8)<1.
\]
There is a graph $H$ on the same entire vertex set, with decidable
adjacency, contained in $G$, retaining every edge from $U$ to $W$ of
weight one, such that for all $A\subseteq U$, $B\subseteq W$ with
$|A|,|B|>hn$,
\[
 \left|e_H(A,B)-\sum_{(a,b)\in A\times B}p(a,b)\right|
 \leq(h/4)|A||B|.
\]
\end{lemma}
\begin{proof}
Index the trials by the oriented existing edges
$I=\{(u,w)\in U\times W:G(u,w)\}$. Disjointness gives a unique
$U$-to-$W$ orientation to each undirected edge meeting both sets.
For $e\in I$, its given weight lies in $[0,1]$, so
\noderef{IndependentUnequalEdgeSampleLaw} supplies the finite product
Bernoulli probability law, with independent coordinate indicators of
expectation $p(e)$ and centered subgaussian parameter $1/4$.
For an outcome, take exactly the selected undirected edges from $I$
on the original vertex set, deleting any edges outside this list.
This graph is simple and contained in $G$; it retains vertices even
if they are isolated. Its finite edge list makes adjacency decidable.

Fix one tested pair $A,B$. Its count is the sum of the indicators
indexed by $I_{AB}=I\cap(A\times B)$, and its expectation is the
sum of their weights. This is the weight sum on all of $A\times B$,
since nonedge weights are zero. Put $r=|I_{AB}|$ and $N=|A||B|$.
We have $r\leq N$ and $N>h^2n^2>0$. For completeness, the
subgaussian conclusion gives the following tail bound without any
equal-probability assumption. Independence yields
\[
 \mathbb E\exp\bigl(s(Z-\mathbb EZ)\bigr)\leq\exp(rs^2/8)
 \quad(s\in\mathbb R).
\]
If $r>0$, Markov's inequality for this nonnegative exponential,
with $s=4t/r$, bounds the upper tail at $t>0$ by
$\exp(-2t^2/r)$. Applying the same argument to the negative centered
sum bounds the lower tail identically. Thus the probability of
$|Z-\mathbb EZ|>t$ is at most $2\exp(-2t^2/r)$, which is at most
$2\exp(-2t^2/N)$ because $r\leq N$. If $r=0$, the discrepancy
is identically zero and the failure probability is zero.
Taking $t=(h/4)N$ therefore bounds the failure probability by
\[
 2\exp(-h^2N/8)\leq2\exp(-h^4n^2/8).
\]
There are $2^n$ subsets of each side, hence $4^n$ subset pairs;
restricting to large pairs does not increase this number. For a
finite family of events, the indicator of the union is bounded by
the sum of their indicators. Integrating gives the finite union
bound, so the probability of any failed test is at most the
left side of the assumed strict cutoff and is strictly less than one.

A trial of weight one succeeds with probability one. The finite
intersection of these probability-one success events still has
probability one, so removing their null complement from the good
test event leaves positive probability. In a finite probability
space a positive-probability event is nonempty. Choose an outcome
in this intersection. Its graph retains all weight-one edges and
satisfies every asserted weak discrepancy bound simultaneously.
\end{proof}
